\documentclass[11pt,a4paper]{article}
\usepackage{gtu-el-yazisi}
\begin{document}
\notbilgi{Elektronik II}{BJT fark yükselteci, CMRR ve aktif yüklü fark yükselteci}{Ders7.pdf}

\sayfa{1}
\okunamadi{In ac}

\elcizimi[0.6]{sekiller/Ders7-s1-1.png}[1]
\begin{align*}
v_1 &= V_{CC} - \frac{\alpha R_C I}{2} - \frac{\alpha\,\Delta i}{2}\footnotemark \\
v_2 &= V_{CC} - \frac{\alpha R_C I}{2} + \frac{\alpha\,\Delta i}{2}
\end{align*}
\footnotetext{Claude notu: muhtemelen $\mp\,\alpha\,\Delta i\,R_C/2$ olmalı ($R_C$ çarpanı eksik; aşağıdaki $v_D = \alpha\,\Delta i\,R_C$ ile uyumlu).}
\formulkutu{v_D = v_2 - v_1 = \alpha\,\Delta i\, R_C}
\[ \Delta i = g_m v_i \]

\elcizimi[0.55]{sekiller/Ders7-s1-2.png}[1]
\[ \Delta i = g_m v_{id} \]
\[ \left.\begin{aligned} v_2 &= \frac{g_m v_{id}}{2}\,R_C \\ v_1 &= -\frac{g_m v_{id}}{2}\,R_C \end{aligned}\right\} \; v_d = g_m v_{id} R_C = v_2 - v_1 \]
\formulkutu{A_d = g_m R_C = \frac{v_d}{v_{id}}}
$A_d \rightarrow$ \textit{differential}

\formulkutu{g_m = \frac{I_C}{V_T}}

\sayfa{2}
\[ A_d = \frac{I_c}{V_T}\,R_C = \frac{I_C R_C}{V_T} \]
soğukta kazanç artar

\section{Effect of $r_o$?}
\elcizimi[0.6]{sekiller/Ders7-s2-1.png}[2]
\elcizimi[0.65]{sekiller/Ders7-s2-2.png}[2]
\[ v_1 = -\frac{v_{id}}{2}\,g_m\,(r_o \parallel R_C) \qquad\qquad v_2 = \frac{v_{id}}{2}\,g_m\,(r_o \parallel R_C) \]
\begin{align*}
v_d = g_m\,(R_C \parallel r_o)\,v_{id} \;\Rightarrow\; A_d = \frac{v_d}{v_{id}} &= g_m\,(R_C \parallel r_o) \\
&\simeq g_m R_C
\end{align*}

\sayfa{3}
\elcizimi[0.6]{sekiller/Ders7-s3-1.png}[3]
\elcizimi[0.45]{sekiller/Ders7-s3-2.png}[3]
\[ R_{id} = 2(\beta + 1)(r_e + R_E) \]
\elcizimi[0.6]{sekiller/Ders7-s3-3.png}[3]

\sayfa{4}
\section{CMRR and Common mode operation}

\elcizimi[0.4]{sekiller/Ders7-s4-1.png}[4]

\section{Ac mode}

\elcizimi[0.65]{sekiller/Ders7-s4-2.png}[4]

\elcizimi[0.55]{sekiller/Ders7-s4-3.png}[4]

\begin{align*}
A_{cm,\,single} = \frac{v_1}{v_{cm}} &\simeq \frac{-R_C}{2R_{EE}+r_e} = \frac{-R_C}{2R_{EE}} \\
\frac{v_2}{v_{cm}} &= \frac{-R_C}{2R_{EE}+r_e} \simeq \frac{-R_C}{2R_{EE}}
\end{align*}
\[ A_{cm,\,diff} = \frac{v_1 - v_2}{v_{cm}} = 0 \]

\sayfa{5}
\section{CMRR diff}

\[ \mathrm{CMRR} = \frac{A_{diff}}{A_{cm\,diff}} \]
\[ \mathrm{CMRR} = \frac{g_m R_C}{0} = \infty \]

\elcizimi[0.5]{sekiller/Ders7-s5-1.png}[5]

\[ \left|\frac{A_{d\,single}}{A_{cm\,single}}\right| = \left|\frac{\frac{1}{2}\,g_m R_C}{\frac{R_C}{2R_{EE}}}\right| = g_m R_{EE} \]

\elcizimi[0.6]{sekiller/Ders7-s5-2.png}[5]

\[ R_{icm} = ? \]
\[ (\beta+1)\Bigl(R_{EE} + \frac{r_e}{2}\Bigr) \]
\[ = (\beta+1)(R_{EE}) \]

\sayfa{6}
\section{Active-Loaded CMOS differential Amplifier}

\elcizimi[0.5]{sekiller/Ders7-s6-1.png}[6]

\elcizimi[0.35]{sekiller/Ders7-s6-2.png}[6]

\elcizimi[0.5]{sekiller/Ders7-s6-3.png}[6]

\elcizimi[0.6]{sekiller/Ders7-s6-4.png}[6]

\elcizimi[0.5]{sekiller/Ders7-s6-5.png}[6]

\[ I_{th} = g_m v_{gs} + \frac{v_{gs}}{r_o} \]
\[ \frac{V_{th}}{I_{th}} = \frac{1}{g_m + 1/r_o} \]
\[ = 1/g_m \]

(Single-ended)

Sadece sağ taraftan çıkış alabiliriz. Peki neden böyle bir şeye ihtiyaç duyuyoruz? Diğer devremiz sadece tek giriş alıyor. \okunamadi{Ona iletiyoruz.}

\sayfa{7}
\elcizimi[0.7]{sekiller/Ders7-s7-1.png}[7]

\elcizimi[0.7]{sekiller/Ders7-s7-2.png}[7]

\begin{align*}
V_{g3} &= -g_{m1}\,\frac{v_{id}}{2}\,\bigl(1/g_{m3} \parallel r_{o3}\bigr) \\
&\simeq -\frac{g_{m1}}{g_{m3}}\,\frac{v_{id}}{2}
\end{align*}
\[ i_o = g_{m4}\,v_{g3} + g_{m2}\,\frac{v_{id}}{2} \]
\begin{align*}
i_o &= \frac{g_{m4}\,g_{m1}}{g_{m3}}\,\frac{v_{id}}{2} + g_m\,\frac{v_{id}}{2} \\
&= g_m\,v_{id}
\end{align*}
\[ g_{m1} = g_{m2},\ g_{m3} = g_{m4} \qquad\qquad = g_{m1,2}\,v_{id} \]

\elcizimi[0.4]{sekiller/Ders7-s7-3.png}[7]

\[ \frac{v_x}{i_x} = R_o \simeq (r_{o2} \parallel r_{o4}) \]
\[ \frac{v_{out}}{v_{id}} = g_m\,(r_{o2} \parallel r_{o4}) \]
\[ r_{o2} = r_{o4} = r_o \]
\[ \Rightarrow\quad = \frac{1}{2}\,g_m r_o \]

\[ A_{cm} \simeq -\frac{1}{2g_{m3}R_{ss}} \]
\[ \mathrm{CMRR} = (g_m r_o)\,g_m R_{ss} \]

\sayfa{8}
\elcizimi[0.45]{sekiller/Ders7-s8-1.png}[8]

p tipinin daha az noise var

\[ G_m = g_{m1} = g_{m2} = g_m \]
\[ R_{out} = r_{o2} \parallel r_{o4} \]
\[ \frac{v_{out}}{v_{id}} = g_m\,(r_{o2} \parallel r_{o4}) \]
\[ r_{o2} = r_{o4} = r_o \]
\[ \frac{v_{out}}{v_{id}} = \frac{1}{2}\,g_m r_o \]

\[ A_{cm} = \frac{-r_{o4}}{\beta_3 R_{EE}} \]
\[ \mathrm{CMRR} = \frac{1}{2}\,\beta_3\,g_m R_{EE} \]

\elcizimi[0.4]{sekiller/Ders7-s8-2.png}[8]

Consider BJT diff amp $I = 0{,}8\,\mathrm{mA}$ \quad $V_A = 100\,\mathrm{V}$, $\beta = 160$. Find $G_m$, $R_o$, $A_d$ and $R_{iod}$. Find $R_{EE}$, $A_{cm}$ and CMRR when the current source is implemented by a single npn current mirror
\[ \left[ A_{cm} = \frac{-r_{o4}}{\beta_3 R_{EE}} \right] \]

\[ g_m = \frac{I_C}{V_T} = ? \qquad V_T \to 25\,\mathrm{mV} \]
\[ g_{m1,2} = \frac{0{,}4\,\mathrm{mA}}{25\,\mathrm{mV}} = 16\,\mathrm{m}\Omega^{-1} \]

\sayfa{9}
\[ R_o = r_{o2} \parallel r_{o4} \]
\[ = \frac{V_A}{I_C} \parallel \frac{V_A}{I_C} = \frac{1}{2}\,\frac{100}{0{,}4} = 125\,\mathrm{k\Omega} \]
\[ A_d = G_m R_o = 16\,\mathrm{m}\Omega^{-1} \times 125\,\mathrm{k\Omega} \]
\[ A_d = 2000\ [\mathrm{V/V}] \]

\[ R_{id} = 2(\beta+1)(r_e) \;\longrightarrow\; = \frac{1}{g_m} \]
\[ = 2 \times 160 \times 62{,}5 \]
\formulkutu{R_{id} = 20\,\mathrm{k\Omega}}

\elcizimi[0.5]{sekiller/Ders7-s9-1.png}[9]

\[ R_{EE} = r_{o5} = \frac{100}{0{,}8} = 125\,\mathrm{k\Omega} \]
\[ |A_{cm}| = \frac{250\,\mathrm{k}}{160\cdot 125\,\mathrm{k\Omega}} = \frac{1}{80} \]

\[ \mathrm{CMRR} = \frac{A_d}{A_{cm}} = \frac{2000}{\frac{1}{80}} = 160\,000 \]
\[ = 20\log 160\,000 \]
\[ = 104{,}08\ \mathrm{dB} \]

\end{document}
